\documentclass[12pt,a5paper]{article}

\usepackage{multicol}
\usepackage{amssymb,amsmath}
\usepackage[utf8]{inputenc}
\usepackage[T2A]{fontenc}
\usepackage[serbianc]{babel}
\usepackage{geometry}
\geometry{a5paper, margin=1.5cm}

% Подешавања заглавља која су смирила компајлер
\addtolength{\topmargin}{-1.59999pt}
\setlength{\headheight}{13.59999pt}

% Модеран и суптилан изглед хедера и футера
\usepackage{fancyhdr}
\pagestyle{fancy}
\fancyhf{}
\fancyhead[L]{\normalfont\small\itshape Математика --- VIII разред}
\fancyhead[R]{\normalfont\small\itshape Дан 4}
\renewcommand{\headrulewidth}{0pt}
\fancyfoot[R]{\normalfont\small Страница \thepage}
\renewcommand{\footrulewidth}{0.2pt}

% Суптилни наслови помоћу titlesec (без масног фонта)
\usepackage{titlesec}
\titleformat{\section}
  {\normalfont\large\scshape}
  {\thesection}
  {1em}
  {}
\titleformat{\subsection}
  {\normalfont\itshape}
  {\thesubsection}
  {1em}
  {}

% Суптилан Садржај помоћу tocloft
\usepackage{tocloft}
\renewcommand{\cftsecfont}{\normalfont\itshape}
\renewcommand{\cftsecpagefont}{\normalfont}
\renewcommand{\cfttoctitlefont}{\normalfont\large\scshape}

\title{Четврти дан: Обла тела}
\author{Припремна настава за поправни испит}
\date{јун 2026. године}

\begin{document}

\maketitle
\thispagestyle{fancy}

\tableofcontents
%\newpage

\section{Ваљак}

Ваљак је обло геометријско тело омеђено са два подударна круга који представљају базе и омотачем који у развијеном облику у равни представља правоугаоник.

\subsection{Формуле и основни задатак}
Површина базе, површина омотача, укупна површина и запремина ваљка рачунају се преко следећих формула у којима константу пи преписујемо:
\begin{align*}
B &= r^2 \pi \\
M &= 2 \cdot r \pi \cdot H \\
P &= 2 \cdot B + M \\
V &= B \cdot H
\end{align*}

Израчунати површину и запремину ваљка ако је задат полупречник основе $r = 5\,\mathrm{cm}$ и висина ваљка $H = 6\,\mathrm{cm}$.
\begin{align*}
B &= 5^2 \pi = 25\pi\,\mathrm{cm}^2 \\
M &= 2 \cdot 5 \pi \cdot 6 = 60\pi\,\mathrm{cm}^2 \\
P &= 2 \cdot 25\pi + 60\pi = 50\pi + 60\pi = 110\pi\,\mathrm{cm}^2 \\
V &= 25\pi \cdot 6 = 150\pi\,\mathrm{cm}^3
\end{align*}

\section{Купа}

Купа је обло геометријско тело које у бази има круг полупречника $r$, док њен омотач чини кружни исечак чија је линија заправо бочна ивица односно изводница купе која се обележава словом $s$.

\subsection{Формуле и основни задатак}
Формуле за израчунавање површине и запремине купе на основном нивоу гласе:
\begin{align*}
B &= r^2 \pi \\
M &= r \cdot \pi \cdot s \\
P &= B + M \\
V &= \frac{1}{3} \cdot B \cdot H
\end{align*}

%\newpage

Израчунати површину купе ако је познат полупречник основе $r = 6\,\mathrm{cm}$ и изводница купе $s = 10\,\mathrm{cm}$.
\begin{align*}
B &= 6^2 \pi = 36\pi\,\mathrm{cm}^2 \\
M &= 6 \cdot \pi \cdot 10 = 60\pi\,\mathrm{cm}^2 \\
P &= 36\pi + 60\pi = 96\pi\,\mathrm{cm}^2
\end{align*}

За израчунавање запремине купе на основном нивоу користе се директно задате вредности за базу и висину. Израчунати запремину купе ако је површина њене базе $B = 45\pi\,\mathrm{cm}^2$, а висина купе износи $H = 8\,\mathrm{cm}$.
\begin{align*}
V &= \frac{1}{3} \cdot B \cdot H \\
V &= \frac{1}{3} \cdot 45\pi \cdot 8 \\
V &= 15\pi \cdot 8 \\
V &= 120\pi\,\mathrm{cm}^3
\end{align*}

\section{Лопта}

Лопта је геометријско тело састављено од свих тарака простора чија је удаљеност од једне фиксиране тачке, центра лопте, мања или једнака задатој дужини која представља полупречник лопте $r$. Лопта нема базу нити омотач, већ је дефинисана јединственом површином.

\subsection{Формуле и основни задатак}
Површина и запремина лопте рачунају се искључиво преко вредности њеног полупречника:
\begin{align*}
P &= 4 \cdot r^2 \pi \\
V &= \frac{4}{3} \cdot r^3 \pi
\end{align*}

Израчунати површину и запремину лопте чији је полупречник $r = 3\,\mathrm{cm}$.
\begin{align*}
P &= 4 \cdot 3^2 \pi = 4 \cdot 9 \pi = 36\pi\,\mathrm{cm}^2 \\
V &= \frac{4}{3} \cdot 3^3 \pi = \frac{4}{3} \cdot 27 \pi = 4 \cdot 9 \pi = 36\pi\,\mathrm{cm}^3
\end{align*}

%\newpage

\section{Задаци за вежбу на часу}

Ученик самостално, уз помоћ формула и смерница наставника, решава следеће задатке ради систематизације знања о облим телима.

\begin{enumerate}
    \item Израчунај површину и запремину ваљка чији је полупречник основе $r = 4\,\mathrm{cm}$, а висина ваљка износи $H = 7\,\mathrm{cm}$.
    \item Одреди површину купе ако је полупречник базе $r = 3\,\mathrm{cm}$, а дужина изводнице износи $s = 5\,\mathrm{cm}$.
    \item Израчунај запремину купе чија је површина основе $B = 27\pi\,\mathrm{cm}^2$, а висина износи $H = 10\,\mathrm{cm}$.
    \item Израчунај површину и запремину лопте чији је полупречник $r = 6\,\mathrm{cm}$.
\end{enumerate}

\end{document}